% ===== PRÉAMBULE CANONIQUE — à coller VERBATIM en tête de CHAQUE .tex =====
\documentclass[11pt,a4paper]{article}
\usepackage[utf8]{inputenc}\usepackage[T1]{fontenc}\usepackage{lmodern}
\usepackage[french]{babel}
\usepackage{amsmath,amssymb,mathtools}\usepackage{icomma}
\usepackage{siunitx}\sisetup{locale=FR}  % \qty{}{}, \si{}, \ang{}
\usepackage{xcolor}\usepackage{colortbl}
\usepackage{tikz}\usepackage{pgfplots}\pgfplotsset{compat=1.18}
\usepackage[siunitx,europeanresistors]{circuitikz} % circuits (élec.)
\usepackage[most]{tcolorbox}\usepackage{enumitem}\usepackage{array,booktabs}
\usepackage[a4paper,margin=2.2cm]{geometry}
\usetikzlibrary{arrows.meta,calc,angles,quotes,positioning,patterns,%
  backgrounds,shapes.geometric,decorations.pathmorphing,decorations.markings,intersections}
\definecolor{primary}{RGB}{222,49,99}\definecolor{primarydark}{RGB}{150,22,55}
\definecolor{defcol}{RGB}{0,114,178}\definecolor{methcol}{RGB}{52,92,148}
\definecolor{excol}{RGB}{0,135,90}\definecolor{attncol}{RGB}{200,40,55}
\tcbset{boxbase/.style={enhanced,breakable,boxrule=0.9pt,arc=2.5pt,
  left=3mm,right=3mm,top=2.3mm,bottom=2.3mm,fonttitle=\bfseries,coltitle=white}}
% --- boîtes TUTORIEL (noms EXACTS, ne pas renommer) ---
\newtcolorbox{cle}[1][]{boxbase,colback=defcol!6,colframe=defcol,
  title={\textsc{À retenir}\ifx\relax#1\relax\relax\else\textmd{ : \itshape #1}\fi}}
\newtcolorbox{methode}[1][]{boxbase,colback=methcol!7,colframe=methcol,sharp corners,
  title={\textsc{Méthode}\ifx\relax#1\relax\relax\else\textmd{ --- \itshape #1}\fi}}
\newtcolorbox{exemplebox}[1][]{boxbase,colback=excol!5,colframe=excol,
  title={\textsc{Exemple}\ifx\relax#1\relax\relax\else\textmd{ : \itshape #1}\fi}}
\newtcolorbox{piege}[1][]{boxbase,colback=attncol!6,colframe=attncol,
  title={\textsc{Piège}\ifx\relax#1\relax\relax\else\textmd{ : \itshape #1}\fi}}
\newtcolorbox{remarque}[1][]{boxbase,colback=black!4,colframe=black!45,
  title={\textsc{Remarque}\ifx\relax#1\relax\relax\else\textmd{ : \itshape #1}\fi}}
% --- boîtes EXERCICE / CORRIGÉ (noms EXACTS, auto-comptées) ---
\newtcolorbox[auto counter]{exo}[1][]{boxbase,colback=primary!4,colframe=primary,
  title={\textsc{Exercice}~\thetcbcounter\ifx\relax#1\relax\relax\else\textmd{ --- \itshape #1}\fi}}
\newtcolorbox[auto counter]{corrige}[1][]{boxbase,colback=excol!4,colframe=excol,
  title={\textsc{Corrigé}~\thetcbcounter\ifx\relax#1\relax\relax\else\textmd{ --- \itshape #1}\fi}}
\newcommand{\vv}[1]{\vec{#1}}\newcommand{\ds}{\displaystyle}
\newcommand{\R}{\mathbb{R}}\newcommand{\N}{\mathbb{N}}\newcommand{\Z}{\mathbb{Z}}
% =========================================================================
\begin{document}
\sisetup{per-mode=symbol}

\begin{corrige}[identifier la force de réaction]
\begin{enumerate}[label=\alph*)]
  \item La force de \emph{gravitation} exercée par le \textbf{livre sur la Terre} (même nature, corps échangés).
  \item La force de \emph{contact} exercée par le \textbf{livre sur la table}. (Attention : la réaction n'est
        pas le poids du livre — celui-ci est exercé par la Terre, pas par la table.)
  \item Elles ont exactement la \textbf{même intensité} (3\textsuperscript{e} loi) ; seules les masses très
        différentes font que la Terre ne « bouge » pas.
\end{enumerate}
\end{corrige}

\begin{corrige}[peut-on frapper plus fort qu'on ne reçoit ?]
\begin{enumerate}[label=\alph*)]
  \item \textbf{Non}, c'est impossible : les deux forces d'une interaction ont \emph{toujours} la même intensité.
        La force de l'élève sur la feuille égale, à chaque instant, celle de la feuille sur l'élève.
  \item Une force n'est pas « possédée » : elle naît de l'interaction. Sur un objet très léger, la force
        de contact reste faible (il « cède » vite) — on ne peut pas lui transmettre plus que ce qu'il renvoie.
\end{enumerate}
\end{corrige}

\begin{corrige}[Rida et William tirent une corde]
\begin{enumerate}[label=\Alph*.]
  \item \textbf{Faux}. Les deux forces aux extrémités de la corde ont la même intensité (3\textsuperscript{e} loi).
  \item \textbf{Faux}. Elles ont la même intensité mais des \emph{sens opposés}.
  \item \textbf{Vrai}. C'est exactement l'énoncé de l'action--réaction.
  \item \textbf{Vrai}. Ce qui diffère, c'est l'interaction avec le \emph{sol} : William prend un meilleur appui
        (force sur le sol plus grande), donc c'est Rida qui est mis en mouvement.
\end{enumerate}
\end{corrige}

\begin{corrige}[la monoroue électrique]
\begin{enumerate}[label=\Alph*.]
  \item \textbf{Vrai}. Action--réaction : même intensité, sens opposés.
  \item \textbf{Vrai}. Chaque force produit une accélération (monoroue vers l'avant, sol/Terre vers l'arrière).
  \item \textbf{Vrai}. Les intensités des deux forces sont égales.
  \item \textbf{Faux}. Les accélérations diffèrent : $a=F/m$, et la masse de la Terre est immense, donc
        son accélération est quasi nulle. (C'est la seule affirmation incorrecte.)
\end{enumerate}
\end{corrige}

\begin{corrige}[le recul du canon, capstone]
\begin{enumerate}[label=\alph*)]
  \item $a_{\text{boulet}}=\dfrac{\Delta v}{\Delta t}=\dfrac{300}{\num{0,02}}=\qty{15000}{\metre\per\second\squared}$ ;
        $F=m_{\text{boulet}}\,a_{\text{boulet}}=6\times15000=\qty{90000}{\newton}$.
  \item Même force sur le canon (réaction) : $F=\qty{90000}{\newton}$, donc
        $a_{\text{canon}}=\dfrac{F}{m_{\text{canon}}}=\dfrac{90000}{1200}=\qty{75}{\metre\per\second\squared}$.
  \item $\dfrac{a_{\text{boulet}}}{a_{\text{canon}}}=\dfrac{15000}{75}=200=\dfrac{1200}{6}=\dfrac{m_{\text{canon}}}{m_{\text{boulet}}}$. \checkmark
\end{enumerate}
\end{corrige}
\end{document}
